% GENERATED FILE. Edit canonical lessons, then run npm run generate:books.
% canonical-source-hash: fnv1a64:87f2dac8f73ab17b
% canonical-lessons: UR-C431-abr, UR-C431-ghirna, UR-C431-chhana, UR-C431-dhund, UR-C431-thamna-let-up

\chapter{Clouds Gather and Clear}
\label{ch:ur-clouds-gather-and-clear}
\hlchaptermodality{431}

\begin{chapteropening}
\textbf{By the end of this chapter:} \emph{I can say the clouds are closing in, the mist is down, and the rain has let up.}

Complete the last lesson of chapter 431: I can say the clouds are closing in, the mist is down, and the rain has let up.
\end{chapteropening}

\section[\texorpdfstring{\ur{ابر} (abr)}{ابر (abr)}]{\ur{ابر} (abr) — a rain cloud; cloud (formal)}

\label{lesson:UR-C431-abr}



\begin{quote}
Before the new one: say the Urdu for airless heat, then the Urdu for to fan (someone).
\end{quote}

\subsection*{You'll want to know: \ur{ابر}}
\textbf{\ur{ابر}} — \emph{abr} — \textquotedblleft{}a rain cloud; cloud (formal)\textquotedblright{}.

\begin{quote}
\ur{آسمان} \ur{پر} \ur{ابر} \ur{ہے۔}
\end{quote}

\begin{quote}
\emph{āsmān par abr hai.}
\end{quote}

\begin{quote}
The sky is overcast.
\end{quote}

Masculine.

\begin{cousinweb}
\textbf{\ur{ابر}} (\emph{abr}) is usually traced to Persian \emph{abr}, \textquotedblleft{}a cloud\textquotedblright{}; it is a more formal word than \textbf{\ur{بادل}}.
\end{cousinweb}

\begin{grammarlens}[title={What you've built}]
One more word for getting around.
\end{grammarlens}

\subsection*{Your turn}
\begin{itemize}
\raggedright
  \item \emph{Say it:} \emph{abr}
  \item \emph{Say it:} \emph{abr}, once more
  \item \emph{Say it:} \emph{āsmān par abr hai.} — The sky is overcast.
  \item \emph{Recall:} say the Urdu for airless heat, then the Urdu for to fan (someone), then say \emph{abr} again
\end{itemize}

\begin{tcolorbox}[breakable,colback=teal!4,colframe=teal!35!black,title={Before you move on}]
What does \ur{ابر} mean? (\textquotedblleft{}A rain cloud; cloud (formal)\textquotedblright{}.) Say it once more.
\end{tcolorbox}
\section[\texorpdfstring{\ur{گھرنا} (ghirnā)}{گھرنا (ghirnā)}]{\ur{گھرنا} (ghirnā) — to be surrounded; to close in (of clouds)}

\label{lesson:UR-C431-ghirna}



\begin{quote}
Before the new one: say the Urdu for to fan (someone), then the Urdu for a rain cloud.
\end{quote}

\subsection*{You'll want to know: \ur{گھرنا}}
\textbf{\ur{گھرنا}} — \emph{ghirnā} — \textquotedblleft{}to be surrounded; to close in (of clouds)\textquotedblright{}.

\begin{quote}
\ur{پہاڑوں} \ur{پر} \ur{بادل} \ur{گھر} \ur{آئے۔}
\end{quote}

\begin{quote}
\emph{pahāṛoṅ par bādal ghir āe.}
\end{quote}

\begin{quote}
Clouds closed in over the mountains.
\end{quote}

Past (it …): \ur{وہ} \ur{آسمان} \ur{بادلوں} \ur{سے} \ur{گھر} \ur{گیا}, \emph{voh āsmān bādaloṅ se ghir gayā}.

\begin{cousinweb}
\textbf{\ur{گھرنا}} (\emph{ghirnā}) is inherited: CDIAL lists the Hindi cousin \emph{ghirnā}, \textquotedblleft{}to be surrounded\textquotedblright{}, under \emph{ghir}, \textquotedblleft{}go round\textquotedblright{}, a form it reconstructs (CDIAL 4474).
\end{cousinweb}

\begin{grammarlens}[title={What you've built}]
One more word for getting around.
\end{grammarlens}

\subsection*{Your turn}
\begin{itemize}
\raggedright
  \item \emph{Say it:} \emph{ghirnā}
  \item \emph{Say it:} \emph{ghirnā}, once more
  \item \emph{Say it:} \emph{pahāṛoṅ par bādal ghir āe.} — Clouds closed in over the mountains.
  \item \emph{Recall:} say the Urdu for to fan (someone), then the Urdu for a rain cloud, then say \emph{ghirnā} again
\end{itemize}

\begin{tcolorbox}[breakable,colback=teal!4,colframe=teal!35!black,title={Before you move on}]
What does \ur{گھرنا} mean? (\textquotedblleft{}To be surrounded; to close in (of clouds)\textquotedblright{}.) Say it once more.
\end{tcolorbox}
\section[\texorpdfstring{\ur{چھانا} (chhānā)}{چھانا (chhānā)}]{\ur{چھانا} (chhānā) — to spread over; to cover (of clouds or mist)}

\label{lesson:UR-C431-chhana}



\begin{quote}
Before the new one: say the Urdu for a rain cloud, then the Urdu for to be surrounded.
\end{quote}

\subsection*{You'll want to know: \ur{چھانا}}
\textbf{\ur{چھانا}} — \emph{chhānā} — \textquotedblleft{}to spread over; to cover (of clouds or mist)\textquotedblright{}.

\begin{quote}
\ur{شام} \ur{کو} \ur{آسمان} \ur{پر} \ur{ابر} \ur{چھا} \ur{گیا۔}
\end{quote}

\begin{quote}
\emph{shām ko āsmān par abr chhā gayā.}
\end{quote}

\begin{quote}
In the evening cloud spread across the sky.
\end{quote}

Past (it …): \ur{وہ} \ur{کہرا} \ur{شہر} \ur{پر} \ur{چھا} \ur{گیا}, \emph{voh kohrā shahr par chhā gayā}.

\begin{cousinweb}
\textbf{\ur{چھانا}} (\emph{chhānā}) is inherited: CDIAL lists the Hindi cousin \emph{chānā}, \textquotedblleft{}to thatch, to spread around\textquotedblright{}, under Sanskrit \emph{chādáyati}, \textquotedblleft{}covers\textquotedblright{} (CDIAL 5018).
\end{cousinweb}

\begin{grammarlens}[title={What you've built}]
One more word for getting around.
\end{grammarlens}

\subsection*{Your turn}
\begin{itemize}
\raggedright
  \item \emph{Say it:} \emph{chhānā}
  \item \emph{Say it:} \emph{chhānā}, once more
  \item \emph{Say it:} \emph{shām ko āsmān par abr chhā gayā.} — In the evening cloud spread across the sky.
  \item \emph{Recall:} say the Urdu for a rain cloud, then the Urdu for to be surrounded, then say \emph{chhānā} again
\end{itemize}

\begin{tcolorbox}[breakable,colback=teal!4,colframe=teal!35!black,title={Before you move on}]
What does \ur{چھانا} mean? (\textquotedblleft{}To spread over; to cover (of clouds or mist)\textquotedblright{}.) Say it once more.
\end{tcolorbox}
\section[\texorpdfstring{\ur{دھند} (dhund)}{دھند (dhund)}]{\ur{دھند} (dhund) — mist; haze}

\label{lesson:UR-C431-dhund}



\begin{quote}
Before the new one: say the Urdu for to be surrounded, then the Urdu for to spread over.
\end{quote}

\subsection*{You'll want to know: \ur{دھند}}
\textbf{\ur{دھند}} — \emph{dhund} — \textquotedblleft{}mist; haze\textquotedblright{}.

\begin{quote}
\ur{صبح} \ur{دریا} \ur{پر} \ur{دھند} \ur{تھی۔}
\end{quote}

\begin{quote}
\emph{subah daryā par dhund thī.}
\end{quote}

\begin{quote}
There was mist on the river in the morning.
\end{quote}

Feminine.

\begin{cousinweb}
\textbf{\ur{دھند}} (\emph{dhund}) is inherited: CDIAL lists the Hindi cousin \emph{dhundh}, \textquotedblleft{}haze\textquotedblright{}, under \emph{dhūmāndha}, \textquotedblleft{}dim-sighted\textquotedblright{}, built from words for smoke and blindness (CDIAL 6858).
\end{cousinweb}

\begin{grammarlens}[title={What you've built}]
One more word for getting around.
\end{grammarlens}

\subsection*{Your turn}
\begin{itemize}
\raggedright
  \item \emph{Say it:} \emph{dhund}
  \item \emph{Say it:} \emph{dhund}, once more
  \item \emph{Say it:} \emph{subah daryā par dhund thī.} — There was mist on the river in the morning.
  \item \emph{Recall:} say the Urdu for to be surrounded, then the Urdu for to spread over, then say \emph{dhund} again
\end{itemize}

\begin{tcolorbox}[breakable,colback=teal!4,colframe=teal!35!black,title={Before you move on}]
What does \ur{دھند} mean? (\textquotedblleft{}Mist; haze\textquotedblright{}.) Say it once more.
\end{tcolorbox}
\section[\texorpdfstring{\ur{تھمنا} (thamnā)}{تھمنا (thamnā)}]{\ur{تھمنا} (thamnā) — to let up; to stop (of rain or wind)}

\label{lesson:UR-C431-thamna-let-up}



\begin{quote}
Before the new one: say the Urdu for to spread over, then the Urdu for mist.
\end{quote}

\subsection*{You'll want to know: \ur{تھمنا}}
\textbf{\ur{تھمنا}} — \emph{thamnā} — \textquotedblleft{}to let up; to stop (of rain or wind)\textquotedblright{}.

\begin{quote}
\ur{دوپہر} \ur{تک} \ur{بارش} \ur{تھم} \ur{گئی۔}
\end{quote}

\begin{quote}
\emph{dopahar tak bārish tham gaī.}
\end{quote}

\begin{quote}
By midday the rain had let up.
\end{quote}

Past (it …): \ur{وہ} \ur{طوفان} \ur{تھم} \ur{گیا}, \emph{voh tūfān tham gayā}.

\begin{cousinweb}
is inherited: CDIAL lists the Hindi cousin \emph{thāmnā}, \textquotedblleft{}to prop, to stop\textquotedblright{}, under Sanskrit \emph{stambhate} (CDIAL 13683); \textbf{\ur{تھمنا}} is the form for something that stops by itself, \textbf{\ur{تھمانا}} the one for handing something over.
\end{cousinweb}

\begin{grammarlens}[title={What you've built}]
One more word for getting around.
\end{grammarlens}

\subsection*{Your turn}
\begin{itemize}
\raggedright
  \item \emph{Say it:} \emph{thamnā}
  \item \emph{Say it:} \emph{thamnā}, once more
  \item \emph{Say it:} \emph{dopahar tak bārish tham gaī.} — By midday the rain had let up.
  \item \emph{Recall:} say the Urdu for to spread over, then the Urdu for mist, then say \emph{thamnā} again
\end{itemize}

\begin{tcolorbox}[breakable,colback=teal!4,colframe=teal!35!black,title={Before you move on}]
What does \ur{تھمنا} mean? (\textquotedblleft{}To let up; to stop (of rain or wind)\textquotedblright{}.) Say it once more.
\end{tcolorbox}
